\section{Related Work}
\label{sec:related-work}

\paragraph{Reverse engineering beyond the web ecosystem.}
Native-code reverse engineering provides closely related techniques for identifying transformed functions and reused libraries. BinHunt~\cite{gao2008binhunt} combines graph matching with symbolic execution and theorem proving to compare basic-block computations, whereas Gemini~\cite{xu2017gemini} learns control-flow-graph embeddings for cross-platform binary-function similarity. Asm2Vec~\cite{ding2019asm2vec} instead learns representations from assembly instructions to support clone search under compiler optimization and obfuscation. For binary-to-source reuse detection, B2SFinder~\cite{feng2019b2sfinder} combines weighted matching of compilation-resistant features with project structure, while BinaryAI~\cite{jiang2024binaryai} retrieves source functions using transformer embeddings and refines matches using link-time locality and function-call relationships. The latter is particularly relevant to ReCall because both combine individual-function similarity with surrounding structural evidence. More recently, the BScout preprint~\cite{li2026bscout} supplements unresolved functions' candidates using the positions of existing matches and employs LLMs to select and review correspondences, also addressing inlined compositions through specialized training. ReCall shares the use of learned similarity and structural context, but targets small library functions in bundled JavaScript: POG anchors initialize exploration over bundler-aware function relations, and candidate bodies are compared under assumed external-binding correspondence. Thus, ReCall's distinction lies in its recovery workflow for bundled JavaScript, rather than in the general use of embeddings, neighboring matches, or LLM judgments for reverse engineering.
